% =============================================================
%  appendice_adt.tex
%  Appendice B --- Abstract Data Type (ADT)
% =============================================================

\chapter{Abstract Data Type (ADT)}
\label{app:adt}

\begin{intuizione}
{\large\bfseries Cos'è un ADT?}

\medskip
Un tipo di dato definito \textbf{solo} dalle operazioni che supporta,\\
\textcolor{black!60}{non da come è implementato internamente.}

\medskip
\textcolor{csTeal}{\itshape Chi usa un ADT vede l'interfaccia.
Chi lo implementa sceglie la rappresentazione.
Le due cose sono separate.}
\end{intuizione}

\vspace{4pt}

% ================================================================
\section{L'idea fondamentale}
% ================================================================

Quando programmi, lavori continuamente con strutture dati: liste,
array, hash table. Ma c'è una distinzione importante tra la
\emph{struttura concreta} che usi e il \emph{concetto astratto}
che rappresenta.

Considera uno stack. Dal punto di vista concettuale, uno stack è
una collezione con una regola precisa: l'ultimo elemento inserito
è il primo a uscire (LIFO). Questa è l'\textbf{essenza} dello stack
--- non importa se sotto c'è una lista, un array, o qualcos'altro.

L'idea dell'ADT è esattamente questa: separare \emph{cosa} fa
una struttura dati da \emph{come} lo fa. Chi usa lo stack vede
solo le operazioni (\fn{push}, \fn{pop}, \fn{peek}). Chi lo implementa
sceglie la rappresentazione interna. Le due prospettive sono
indipendenti.

\begin{concetto}{I tre elementi di un ADT}
\keyline{Interfaccia}{Le operazioni visibili al chiamante: \fn{push}, \fn{pop}, \fn{empty?}.}
\keyline{Implementazione}{La struttura dati sottostante: lista, array, hash table.}
\keyline{Invariante}{Proprietà che l'ADT garantisce sempre (es: LIFO per lo stack).}
\end{concetto}

Questa separazione ha una conseguenza pratica importante:
puoi cambiare l'implementazione senza modificare il codice che
usa l'ADT. Se oggi il tuo stack è una lista e domani diventa
un array (magari per motivi di performance), il resto del programma
continua a funzionare esattamente come prima.

\begin{esempio}[Stack: stessa interfaccia, implementazioni diverse]
\begin{lstlisting}
;; Interfaccia dello stack (quello che il chiamante vede)
;; - make-stack    : crea stack vuoto
;; - push item s   : aggiunge in cima
;; - pop s         : rimuove e restituisce la cima
;; - peek s        : legge la cima senza rimuovere
;; - empty? s      : vero se vuoto

;; Implementazione A: lista semplice
(defun make-stack-A () nil)
(defun push-A (item s) (cons item s))
(defun pop-A (s) (values (car s) (cdr s)))
(defun peek-A (s) (car s))
(defun empty-A? (s) (null s))

;; Implementazione B: vettore con indice
(defun make-stack-B ()
  (list (make-array 100 :fill-pointer 0)))
(defun push-B (item s)
  (vector-push item (car s)) s)
(defun pop-B (s)
  (vector-pop (car s)))
\end{lstlisting}

\detail{Il codice che usa lo stack non cambia se passi da A a B.
Vede solo \fn{push}/\fn{pop}, non sa cosa c'è sotto.}
\end{esempio}

% ================================================================
\section{Perché usare gli ADT?}
% ================================================================

\begin{concetto}{I tre vantaggi degli ADT}
\keyline{Astrazione}{Nascondi la complessità. Il chiamante pensa a ``stack'', non a ``lista con cons in testa''.}
\keyline{Modularità}{Puoi cambiare implementazione senza toccare il codice che usa l'ADT.}
\keyline{Correttezza}{L'interfaccia forza l'uso corretto. Non puoi accidentalmente rompere l'invariante.}
\end{concetto}

\begin{attenzione}
\textbf{Senza ADT: accesso diretto alla rappresentazione.}

\begin{lstlisting}
;; Stack "nudo" come lista
(defparameter *stack* '(3 2 1))

;; Chiunque puo' fare cose sbagliate:
(setf *stack* (append *stack* '(4)))  ; aggiunge in fondo! non e' LIFO
(setf (cadr *stack*) 99)              ; modifica elemento interno!
\end{lstlisting}

Con un ADT, queste operazioni non sarebbero possibili:
l'interfaccia espone solo \fn{push}/\fn{pop}.
\end{attenzione}

% ================================================================
\section{ADT in Common Lisp}
% ================================================================

Common Lisp non ha un costrutto ``ADT'' esplicito come alcuni linguaggi
(ML, Haskell). Ma puoi ottenere lo stesso risultato con:

\begin{concetto}{Tre modi per implementare ADT in Lisp}
\keyline{\fn{defstruct}}{Campi privati, accessor generati. Semplice e veloce.}
\keyline{Closure}{Funzioni che catturano stato privato. Massimo incapsulamento.}
\keyline{CLOS}{Classi con slot e metodi. Massima flessibilità.}
\end{concetto}

\begin{esempio}[Stack con closure: stato completamente nascosto]
\begin{lstlisting}
(defun make-stack ()
  (let ((data '()))   ; stato privato, invisibile dall'esterno
    (lambda (op &optional item)
      (case op
        (:push  (push item data))
        (:pop   (pop data))
        (:peek  (car data))
        (:empty (null data))))))

;; Uso
(defparameter *s* (make-stack))
(funcall *s* :push 1)
(funcall *s* :push 2)
(funcall *s* :pop)     ; => 2
(funcall *s* :peek)    ; => 1

;; Non c'e' modo di accedere a 'data' direttamente
\end{lstlisting}

\detail{La variabile \fn{data} è catturata dalla closure.
L'unico modo di interagire è attraverso i messaggi
\fn{:push}, \fn{:pop}, \fn{:peek}, \fn{:empty}.}
\end{esempio}

% ================================================================
\section{ADT classici}
% ================================================================

\begin{tcolorbox}[
  enhanced,
  colback=black!2, colframe=black!20,
  arc=3pt, boxrule=0.4pt,
  left=12pt, right=12pt, top=9pt, bottom=9pt,
  before skip=8pt, after skip=10pt
]
\small

\begin{tabular}{@{}llp{5.5cm}@{}}
\toprule
\textbf{ADT} & \textbf{Operazioni chiave} & \textbf{Invariante} \\
\midrule
Stack & \fn{push}, \fn{pop}, \fn{peek} & LIFO: ultimo inserito = primo estratto \\[4pt]
Queue & \fn{enqueue}, \fn{dequeue} & FIFO: primo inserito = primo estratto \\[4pt]
Set & \fn{add}, \fn{remove}, \fn{member?} & No duplicati \\[4pt]
Map/Dict & \fn{get}, \fn{put}, \fn{remove} & Chiavi uniche \\[4pt]
Priority Queue & \fn{insert}, \fn{extract-min} & Estrae sempre il minimo \\[4pt]
BST & \fn{insert}, \fn{find}, \fn{delete} & Ordinamento: sx $<$ nodo $<$ dx \\
\bottomrule
\end{tabular}
\end{tcolorbox}

\detail{Ogni ADT può avere implementazioni diverse con trade-off diversi.
Una coda può essere implementata con due liste (efficiente per \fn{enqueue}/\fn{dequeue})
o con un array circolare (efficiente in spazio).
L'interfaccia resta la stessa.}

% ================================================================
\section{Quando serve un ADT?}
% ================================================================

\begin{intuizione}
{\bfseries Non sempre.}

\medskip
Se usi una lista come stack in tre righe di codice,
incapsularla in un ADT è overhead inutile.

\medskip
\textcolor{csTeal}{\itshape Un ADT serve quando:}
\begin{itemize}
  \item La struttura è usata in più punti del programma
  \item Potresti voler cambiare implementazione in futuro
  \item L'invariante è importante e facile da violare
  \item Vuoi documentare l'intento (``questo è uno stack'')
\end{itemize}

\medskip
\textcolor{black!60}{\small Regola pratica: inizia semplice.
Se la rappresentazione ``nuda'' crea problemi, incapsula.}
\end{intuizione}

